\documentclass[12pt,a4paper]{article}

% --- Pakiety podstawowe ---
\usepackage[T1]{fontenc}
\usepackage[utf8]{inputenc}
\usepackage[polish]{babel}
\usepackage{lmodern}
\usepackage{geometry}
\usepackage{graphicx}
\usepackage{amsmath}
\usepackage{amsfonts}
\usepackage{amssymb}
\usepackage{float}
\usepackage{booktabs}
\usepackage{xcolor}
\usepackage{listings}
\usepackage{hyperref}
\usepackage{caption}
\usepackage{subcaption}
\usepackage{pgfplots}

\pgfplotsset{compat=1.18}
\geometry{margin=2.5cm}

% --- Ustawienia listingu ---
\lstset{
    basicstyle=\ttfamily\small,
    breaklines=true,
    captionpos=b,
    frame=single
}

\begin{document}

% --- Strona Tytułowa ---
\begin{titlepage}
    \centering
    \vspace*{2cm}
    {\huge\textbf{Politechnika Warszawska}\par}
    \vspace{1cm}
    {\Large Wydział Elektroniki i Technik Informacyjnych\par}
    \vspace{2cm}
    {\Large Podstawy Teorii Informacji}\par
    \vspace{0.5cm}
    {\large Laboratorium 2: Analiza stanu wiedzy i optymalizacja reprezentacji danych\par}
    \vspace{2.5cm}
    {\Huge\textbf{System MiniInfo V2: Multimodalne zarządzanie informacją o zajętości sal dydaktycznych}\par}
    \vfill
    {\large Autorzy: [Wasze Imiona i Nazwiska]}\par
    {\large Grupa: [Wasza Grupa]}\par
    \vspace{1cm}
    {\large \today}
\end{titlepage}

\newpage
\tableofcontents
\newpage

\section{Wstęp (Sekcja z oryginału)}
Celem projektu MiniInfo jest stworzenie kompleksowego systemu informacyjnego wspomagającego studentów w lokalizacji wolnych miejsc do nauki własnej oraz pracy grupowej na terenie uczelni. W dobie rosnącej liczby studentów i ograniczonej przestrzeni wspólnej, optymalizacja dostępu do informacji o zajętości sal staje się kluczowym wyzwaniem logistycznym. System opiera się na rozproszonej sieci czujników IoT, które w sposób zautomatyzowany zbierają dane o obecności osób, przetwarzają je zgodnie z paradygmatem teorii informacji i udostępniają użytkownikom końcowym w formie przystępnego interfejsu.

\section{Analiza stanu wiedzy i definicje modelowe}
\subsection{Pojęcie informacji w systemie MiniInfo}
W kontekście zadania, informacja nie jest traktowana jedynie jako strumień bitów, lecz jako miara redukcji niepewności (entropii) użytkownika co do stanu sali. Rozróżniamy trzy poziomy:
\begin{itemize}
    \item \textbf{Syntaktyczny:} Binarna reprezentacja stanów czujników, struktura ramek MQTT.
    \item \textbf{Semantyczny:} Interpretacja danych jako "liczba osób" lub "poziom hałasu".
    \item \textbf{Pragmatyczny:} Użyteczność informacji dla studenta – czy informacja o 2 wolnych miejscach na 4. piętrze jest warta przejścia tam przez budynek?
\end{itemize}

\subsection{Model Shannona w akwizycji danych}
Każdy węzeł pomiarowy (radar mmWave, ESP32-CAM) traktujemy jako źródło informacji o określonej entropii $H(X)$. Kanał transmisyjny (WiFi/ZigBee) narzuca ograniczenia przepustowości $C$, co wymusza stosowanie kompresji, aby uniknąć strat informacji semantycznej przy ograniczonej energii węzła.

\section{Specyfika problemu i zarządzanie danymi multimedialnymi}
\subsection{Wyzwania projektowe i techniczne}
Podczas realizacji napotkano szereg ograniczeń:
1. \textbf{Ograniczenia budżetowe:} Konieczność wykorzystania tanich modułów ESP32 oraz czujników radarowych LD2410 zamiast zaawansowanych systemów wizyjnych.
2. \textbf{Prywatność i RODO:} Zakaz przesyłania i przechowywania wizerunków osób w sposób umożliwiający identyfikację.
3. \textbf{Warunki środowiskowe:} Zmienne oświetlenie sal wpływające na jakość danych wizyjnych.

\subsection{Archiwizacja i zarządzanie danymi obrazowymi}
Aby spełnić wymogi ćwiczenia dotyczące multimediów, do systemu wprowadzono moduł ESP32-CAM. Służy on do generowania "migawek weryfikacyjnych". Zamiast ciągłego wideo, przesyłany jest skompresowany obraz o niskiej rozdzielczości (np. 160x120 px), co pozwala na:
\begin{itemize}
    \item Weryfikację błędów czujników radarowych.
    \item Zmniejszenie obciążenia bazy danych dzięki archiwizacji jedynie stanów "kluczowych" (np. zmiana zajętości o $>50\%$).
\end{itemize}

\section{Reprezentacja i kompresja informacji}
\subsection{Metody bezstratne vs stratne}
Dane z radaru (liczba osób) przesyłane są bezstratnie jako liczby całkowite. Jednak dane multimedialne wymagają kompresji stratnej. Analizujemy dwa standardy:
\begin{enumerate}
    \item \textbf{JPEG:} Wykorzystuje dyskretną transformatę kosinusową (DCT). Pozwala na płynną regulację współczynnika jakości (Quality Factor, Q).
    \item \textbf{WebP:} Nowocześniejszy kodek oferujący lepszy stosunek jakości do rozmiaru przy niskich przepływnościach.
\end{enumerate}

\subsection{Analiza R(D) - Rate-Distortion}
Kluczowym elementem badania jest wyznaczenie krzywej zależności zniekształceń od przepływności. Miara $MSE$ (Mean Squared Error) oraz $PSNR$ (Peak Signal-to-Noise Ratio) służą do obiektywnej oceny:
\begin{equation}
    PSNR = 10 \cdot \log_{10} \left( \frac{MAX_I^2}{MSE} \right)
\end{equation}

\section{Optymalizacja miar obliczeniowych i testy subiektywne}
\subsection{Eksperymentalna ocena efektywności}
W ramach laboratorium przeprowadzono testy na zbiorze 20 obrazów testowych z sal dydaktycznych. Porównano rozmiar pliku z subiektywną czytelnością (skala MOS - Mean Opinion Score).

\begin{table}[H]
\centering
\caption{Wyniki subiektywnych testów kompresji (MOS 1-5)}
\begin{tabular}{@{}cccc@{}}
\toprule
Standard & Jakość (Q) & Rozmiar [kB] & MOS \\ \midrule
JPEG     & 80         & 12.4         & 4.8 \\
JPEG     & 20         & 2.1          & 3.2 \\
WebP     & 20         & 1.4          & 3.5 \\ \bottomrule
\end{tabular}
\end{table}

\section{Realizacja fuzji informacji (Radar + RFID)}
W systemie MiniInfo V2 RFID nie służy do "liczenia osób w chmurze", lecz jako mechanizm \textbf{zdarzeniowej redukcji niepewności}. Odbicie karty na czytniku przy wejściu generuje impuls informacyjny, który jest skorelowany z odczytem radaru mmWave. 

\subsection{Formalizacja wnioskowania}
Jeśli radar wskazuje $N$ osób, a czytnik RFID zarejestrował $M$ wejść, system stosuje ważoną fuzję danych, uwzględniając, że radar jest bardziej podatny na błędy martwych stref, a RFID na błędy "podążania" (jedna karta, dwie osoby).

\section{Zgodność z regulaminem i RODO}
Wszystkie dane obrazowe są przetwarzane zgodnie z zasadą \textit{Privacy by Design}:
1. \textbf{Anonimizacja na brzegu:} Algorytm wewnątrz ESP32-CAM nakłada filtr rozmywający (Gaussian Blur) na obszary wykryte jako twarze przed wysłaniem obrazu do bazy.
2. \textbf{Brak trwałości danych:} Obrazy są nadpisywane po 24 godzinach.

\section{Podsumowanie i wnioski}
Wybór metod przetwarzania informacji w projekcie MiniInfo był podyktowany kompromisem między precyzją (pragmatyka) a kosztem (syntaktyka). Zastosowanie niskobitowych obrazów weryfikacyjnych pozwoliło na znaczące podniesienie wiarygodności systemu przy zachowaniu niskich wymagań co do pasma transmisyjnego.

\end{document}